% Bagian: sets-functions-relations
% Bab: sets
% Subbagian: reduction

\documentclass[../../../include/open-logic-section]{subfiles}

\begin{document}

\olfileid[id]{sfr}{siz}{red}

\olsection{Reduksi}

\begin{editorial}
  Bagian ini membuktikan ketakterhitungan dengan reduksi, selaras dengan
  hasil-hasil dalam \olref[nen]{sec}. Versi alternatif yang sedikit lebih
  ringkas dan selaras dengan hasil-hasil dalam \olref[nen-alt]{sec} diberikan
  dalam \olref[red-alt]{sec}.
\end{editorial}

Kita telah menunjukkan bahwa $\Pow{\PosInt}$ !!{nonenumerable} dengan argumen
diagonalisasi. Kita juga telah mempunyai bukti bahwa $\Bin^\omega$, himpunan
semua barisan tak hingga dari $0$ dan $1$, adalah !!{nonenumerable}. Berikut
cara lain untuk membuktikan bahwa $\Pow{\PosInt}$ !!{nonenumerable}:
Tunjukkan bahwa \emph{jika $\Pow{\PosInt}$ !!{enumerable}, maka $\Bin^\omega$
juga !!{enumerable}}. Karena kita tahu bahwa $\Bin^\omega$ tidak
!!{enumerable}, $\Pow{\PosInt}$ juga tidak mungkin demikian. Ini disebut
\emph{mereduksi} satu masalah ke masalah lain---dalam kasus ini, kita
mereduksi masalah mengenumerasi $\Bin^\omega$ ke masalah mengenumerasi
$\Pow{\PosInt}$. Penyelesaian bagi masalah yang kedua---suatu enumerasi
$\Pow{\PosInt}$---akan menghasilkan penyelesaian bagi masalah yang
pertama---suatu enumerasi $\Bin^\omega$.

Bagaimana kita mereduksi masalah mengenumerasi suatu himpunan~$B$ ke masalah
mengenumerasi suatu himpunan~$A$? Kita memberikan cara untuk mengubah suatu
enumerasi~$A$ menjadi enumerasi~$B$. Cara termudah untuk melakukannya ialah
mendefinisikan fungsi !!a{surjective} $f\colon A \to B$. Jika $x_1$, $x_2$,
\dots{} mengenumerasi~$A$, maka $f(x_1)$, $f(x_2)$, \dots{} akan mengenumerasi
~$B$. Dalam kasus kita, kita mencari fungsi surjektif $f\colon
\Pow{\PosInt} \to \Bin^\omega$.

\begin{prob}
Tunjukkan bahwa jika terdapat fungsi !!{injective} $g\colon B \to A$ dan
$B$~adalah !!{nonenumerable}, maka $A$ juga demikian. Lakukan ini
dengan menunjukkan cara menggunakan~$g$ untuk mengubah suatu enumerasi~$A$
menjadi enumerasi~$B$.
\end{prob}

\begin{proof}[Bukti reduksi untuk {\olref[nen]{thm:nonenum-pownat}}]
Andaikan $\Pow{\PosInt}$ adalah !!{enumerable}, sehingga terdapat suatu
enumerasi darinya, yaitu $Z_{1}$, $Z_{2}$, $Z_{3}$, \dots

Definisikan fungsi $f \colon \Pow{\PosInt} \to \Bin^\omega$ dengan menetapkan
$f(Z)$ sebagai barisan $s$ sedemikian sehingga $s(n) = 1$ jika dan hanya jika
$n \in Z$, dan $s(n) = 0$ jika tidak. Ini jelas mendefinisikan
suatu fungsi, karena jika $Z \subseteq \PosInt$, maka setiap $n \in \PosInt$
merupakan !!a{element} dari $Z$ atau bukan. Sebagai contoh, himpunan bilangan
genap positif $2\PosInt = \{2, 4, 6, \dots\}$ dipetakan ke barisan
$010101\dots$, himpunan kosong dipetakan ke $0000\dots$, dan himpunan
$\PosInt$ sendiri dipetakan ke $1111\dots$.

Fungsi itu juga !!{surjective}: Setiap barisan $0$ dan $1$ bersesuaian dengan
suatu himpunan bilangan bulat positif, yakni himpunan yang beranggotakan
bilangan-bilangan bulat yang bersesuaian dengan posisi tempat barisan tersebut
bernilai~$1$. Secara lebih saksama, andaikan $s \in \Bin^\omega$. Definisikan
$Z \subseteq \PosInt$ sebagai:
\[
Z = \Setabs{n \in \PosInt}{s(n) = 1}
\]
Maka $f(Z) = s$, sebagaimana dapat diperiksa dari definisi~$f$.

Sekarang perhatikan daftar
\[
f(Z_1), f(Z_2), f(Z_3), \dots
\]
Karena $f$ !!{surjective}, setiap anggota $\Bin^\omega$ harus muncul sebagai
nilai~$f$ untuk suatu argumen, sehingga harus muncul dalam daftar. Jadi,
daftar ini harus mengenumerasi seluruh~$\Bin^\omega$.

Dengan demikian, jika $\Pow{\PosInt}$ !!{enumerable}, maka $\Bin^\omega$ juga
!!{enumerable}. Namun, $\Bin^\omega$ adalah !!{nonenumerable}
(\olref[nen]{thm:nonenum-bin-omega}). Jadi, $\Pow{\PosInt}$ adalah
!!{nonenumerable}.
\end{proof}

\begin{explain}
Kita mudah keliru mengenai arah reduksi. Sebagai contoh, fungsi
!!a{surjective} $g \colon \Bin^\omega \to B$ \emph{tidak} membuktikan bahwa
$B$ !!{nonenumerable}. (Perhatikan $g \colon \Bin^\omega \to \Bin$ yang
didefinisikan oleh $g(s) = s(1)$, yakni fungsi yang memetakan suatu barisan
$0$ dan $1$ ke !!{element} pertamanya. Fungsi itu !!{surjective}, karena
beberapa barisan dimulai dengan $0$ dan beberapa lainnya dengan $1$. Namun,
$\Bin$ hingga.) Perhatikan pula bahwa fungsi~$f$ harus !!{surjective}; jika
tidak, argumennya tidak berlaku: $f(x_1)$, $f(x_2)$, \dots{} tidak dijamin
memuat semua !!{element}s~$B$. Sebagai contoh,
\[
h(n) = \underbrace{000\dots0}_{\text{$n$ buah $0$}}
\]
mendefinisikan fungsi $h\colon \PosInt \to \Bin^\omega$, tetapi $\PosInt$
adalah !!{enumerable}.
\end{explain}

\begin{prob}
Tunjukkan dengan argumen reduksi bahwa himpunan dari semua \emph{himpunan
pasangan} bilangan bulat positif adalah !!{nonenumerable}.
\end{prob}

\begin{prob}\label{sfr:siz:red:prob:nat-nat}
  Tunjukkan dengan argumen reduksi bahwa himpunan~$X$ dari semua fungsi
  $f\colon \Nat \to \Nat$ adalah !!{nonenumerable} (Petunjuk: berikan fungsi
  surjektif dari $X$ ke~$\Bin^\omega$.)
\end{prob}

\begin{prob}
Tunjukkan dengan argumen reduksi bahwa $\Nat^\omega$, himpunan barisan tak
hingga bilangan asli, adalah !!{nonenumerable}.
\end{prob}

\begin{prob}
Misalkan $P$ adalah himpunan fungsi dari himpunan bilangan bulat positif ke
himpunan $\{0\}$, dan misalkan $Q$ adalah himpunan fungsi \emph{parsial} dari
himpunan bilangan bulat positif ke himpunan $\{0\}$. Tunjukkan bahwa $P$~adalah
!!{enumerable} dan $Q$~tidak. (Petunjuk: reduksi masalah mengenumerasi
$\Bin^\omega$ ke masalah mengenumerasi~$Q$.)
\end{prob}

\begin{prob}
Misalkan $S$ adalah himpunan semua fungsi !!{surjective} dari himpunan bilangan
bulat positif ke himpunan \{0,1\}, yakni $S$ terdiri atas semua fungsi
!!{surjective}~$f\colon \PosInt \to \Bin$. Tunjukkan bahwa $S$ adalah
!!{nonenumerable}.
\end{prob}

\begin{prob}
Tunjukkan bahwa himpunan~$\Real$ dari semua bilangan real adalah
!!{nonenumerable}.
\end{prob}

\end{document}
